% Proof-bearing staged module. Namespace nca:. No canonical edits.
\section{The arithmetic of the retained Albert cubic on both lattices}
\label{nca:sec:cubic-arithmetic}

\subsection{Original data and the exact polynomial}

Retain the ordered coordinates \(X=(x_{i,j};d_k)\),
\(1\leq i\leq4\), \(0\leq j\leq5\), \(1\leq k\leq4\), with
\(\sum_{j=0}^5x_{i,j}=0\), and their original Euclidean metric.
Let
\[
 A_5=\{x\in\mathbb Z^6:\sum_jx_j=0\},\qquad
 D_4=\{d\in\mathbb Z^4:\sum_kd_k\equiv0\pmod2\},\qquad
 R=A_5^{\perp4}\perp D_4.
\]
Write \(\lambda=(5,-1,-1,-1,-1,-1)/6\),
\(v_D=(1,0,0,0)\), and \(s_D=(1,1,1,1)/2\).
The original discriminant glue is generated by
\[
 \begin{split}
 &(3,3,3,3;00),\quad(3,3,0,0;10),\quad(3,0,3,0;01),\\
 &(2,2,2,0;00),\quad(2,4,0,2;00)
 \end{split}
\]
in \((\mathbb Z/6)^4\oplus\mathbb F_2^2\), with the same marking
\(10=v_D+D_4\), \(01=s_D+D_4\). Its order is 72: the first
three generators have independent two-primary reductions, and the
last two independent three-primary reductions.  The lattice \(N\)
is its inverse image in \(R^\vee\).

Retain
\[
 \rho_A=(5,3,1,-1,-3,-5)/2,\qquad
 \rho_D=(3,2,1,0),\qquad \rho=\rho_A^{\oplus4}\perp\rho_D,
 \qquad \alpha=(e_0-e_1,0,0,0;0),
\]
and the exact neighbour
\[
 K=\{X\in N:(\rho,X)\equiv0\pmod6\},\qquad
 v=\rho/6-\alpha,\qquad \Lambda=K+\mathbb Zv.
\]
Here \(\rho\in N\), \(\rho^2=84\), and \((\rho,\alpha)=1\),
as follows immediately from these coordinates and the first glue
generator.  The name \(\Lambda\) always denotes this neighbour.

Fix an ordered orthonormal octonion basis \(u_0=1,u_1,\ldots,u_7\).
For the first part of the calculation no multiplication table is
imposed beyond the actual multiplication of these basis elements.
Define, in the original coordinates,
\begin{align}
 s_i&=\sum_{j=0}^5(-1)^jx_{i,j},\label{nca:eq:original-s}\\
 A_{2(i-1)+p}
  &=x_{i,p}-\frac{x_{i,p}+x_{i,p+2}+x_{i,p+4}}3,
       &&p=0,1,\label{nca:eq:original-A}\\
 B_{2(i-1)+p}
  &=x_{i,p+2}-\frac{x_{i,p}+x_{i,p+2}+x_{i,p+4}}3,
       &&p=0,1.\label{nca:eq:original-B}
\end{align}
Write \(A=A_0+\mathbf a\), \(B=B_0+\mathbf b\), where
\(\mathbf a=\sum_{h=1}^7A_hu_h\) and
\(\mathbf b=\sum_{h=1}^7B_hu_h\).  The common mean octonion is
\begin{equation}
 M=\frac{\sqrt2}{6}\sum_{a=0}^3s_{a+1}u_a
       +\frac{\sqrt3}{3}\sum_{k=1}^4d_ku_{k+3}.
 \label{nca:eq:mean}
\end{equation}
The retained isometry therefore has precisely the ordered form
\[
 \mathcal LX=(A+M,B+M,-A-B+M).
\]
This is an identity with its original three triplets: the residuals
sum to zero, the first four common mean coefficients are
\(s_i/\sqrt{18}=\sqrt2s_i/6\), and the last four are
\(d_k/\sqrt3=\sqrt3d_k/3\).
For completeness the original squared norm and the output norm are both
\[
 |A|^2+|B|^2+|A+B|^2+\frac16\sum_i s_i^2+\sum_kd_k^2.
\]
In the original fibres this follows by separating the means
\(s_i/6,-s_i/6\) from their zero-sum triplets; in the output it
follows from \(A+B+(-A-B)=0\) and
\(3|M|^2=\sum_i s_i^2/6+\sum_kd_k^2\). The inverse recovers
\(M=(q_1+q_2+q_3)/3\), then \(s_i=3\sqrt2M_{i-1}\),
\(d_k=\sqrt3M_{k+3}\), and \(A=q_1-M,B=q_2-M\).
Adding \((-1)^ps_i/6\) to the \(p\)-th residual triplet in the
\(i\)-th fibre recovers each original \(x_{i,p+2j}\). Thus the
map is onto and preserves the original pairing by polarization.

For imaginary octonions put
\[
 \mathbf a\times\mathbf b=\frac{\mathbf a\mathbf b-
                                    \mathbf b\mathbf a}{2},\qquad
 \phi(\mathbf a,\mathbf b,\mathbf c)
                     =\langle\mathbf a\times\mathbf b,\mathbf c\rangle.
\]
Let \(\phi_{ijk}=\phi(u_i,u_j,u_k)\).  Thus \(\phi\) is the
actual alternating octonion three-form in the chosen frame, including
all signs. Define
\begin{align}
 D&=A_0^2+A_0B_0+B_0^2-\langle\mathbf a,\mathbf a\rangle
             -\langle\mathbf a,\mathbf b\rangle
             -\langle\mathbf b,\mathbf b\rangle,\label{nca:eq:D}\\
 Z&=(2A_0+B_0)\mathbf a+(A_0+2B_0)\mathbf b
                               -3\mathbf a\times\mathbf b
       =\sum_{h=1}^7Z_hu_h.\label{nca:eq:Z}
\end{align}
Equivalently, the component in the final formula is the exact
quadratic polynomial
\[
 Z_h=(2A_0+B_0)A_h+(A_0+2B_0)B_h
                  -3\sum_{i,j=1}^7\phi_{ijh}A_iB_j.
\]

\begin{theorem}[Exact pullback with every radical and coefficient]
\label{nca:thm:polynomial}
The original cubic
\(F(X)=\operatorname{Re}(((\mathcal LX)_1(\mathcal LX)_2)
                                           (\mathcal LX)_3)\)
has the complete polynomial expression
\begin{equation}
 F(X)=P_0(X)+\sqrt2P_2(X)+\sqrt3P_3(X),
 \label{nca:eq:polynomial}
\end{equation}
where every variable is the displayed linear function of the original
\(x_{i,j},d_k\), and
\begin{align}
 P_0={}&-A_0B_0(A_0+B_0)+B_0\langle\mathbf a,\mathbf a\rangle
              +A_0\langle\mathbf b,\mathbf b\rangle
              +2(A_0+B_0)\langle\mathbf a,\mathbf b\rangle,
                \label{nca:eq:P0}\\
 P_2={}&\frac{-s_1D+s_2Z_1+s_3Z_2+s_4Z_3}{6}
       +\frac{s_1(4s_1^2-3\sum_{i=1}^4s_i^2)}{108}
       -\frac{s_1\sum_{k=1}^4d_k^2}{6},\label{nca:eq:P2}\\
 P_3={}&\frac{d_1Z_4+d_2Z_5+d_3Z_6+d_4Z_7}{3}.
                  \label{nca:eq:P3}
\end{align}
These identities hold throughout the original real vector space, in
particular on \(N\), \(K\), and \(\Lambda\). They retain the ordered
Albert determinant as
\[
 \det Q_{\mathbb O}(\lambda;\mathcal LX)
                =\lambda^3-\lambda(X,X)+2F(X).
\]
\end{theorem}

\begin{proof}
For \(a=a_0+\mathbf a\), \(b=b_0+\mathbf b\),
\(c=c_0+\mathbf c\), binary octonion multiplication gives
\[
 \operatorname{Re}((ab)c)
  =a_0b_0c_0-a_0\langle\mathbf b,\mathbf c\rangle
   -b_0\langle\mathbf a,\mathbf c\rangle
   -c_0\langle\mathbf a,\mathbf b\rangle
   -\phi(\mathbf a,\mathbf b,\mathbf c).
\]
Call the first four, symmetric, terms \(S(a,b,c)\).  Substitute
\((a,b,c)=(A+M,B+M,-A-B+M)\).  Terms with exactly two copies
of \(M\) in \(S\) sum to \(S(M,M,A+B-A-B)=0\). Terms
with exactly one copy sum to
\(-S(M,A,A)-S(M,A,B)-S(M,B,B)\).  The term with no copy is
\(S(A,B,-A-B)=P_0\), by multiplying its four displayed terms.
The alternating term with no copy vanishes; terms with two or
three copies vanish. Its three terms with one copy all equal
\(\phi(\mathbf a,\mathbf b,\operatorname{Im}M)\), since
\[
 \phi(\mathbf a,\operatorname{Im}M,-\mathbf a-\mathbf b)
 =\phi(\mathbf a,\mathbf b,\operatorname{Im}M),\quad
 \phi(\operatorname{Im}M,\mathbf b,-\mathbf a-\mathbf b)
 =\phi(\mathbf a,\mathbf b,\operatorname{Im}M).
\]
Thus all one-copy terms together are
\(-M_0D+\langle\operatorname{Im}M,Z\rangle\).
Finally the term with three copies is
\[
 \operatorname{Re}(M^3)=4M_0^3-3M_0|M|^2,
 \qquad |M|^2=\frac{\sum_i s_i^2}{18}
                              +\frac{\sum_kd_k^2}{3},
 \qquad M_0=\frac{\sqrt2s_1}{6}.
\]
Substitution proves \eqref{nca:eq:P0}--\eqref{nca:eq:P3} with
their displayed factors. In particular it proves the absence of a
\(\sqrt6\) term without discarding any mixed term. The determinant
formula uses the original isometry and its unchanged squared norm.
\end{proof}

\subsection{Finite generators of both original lattices}

Order the 24 original simple roots as
\[
 g_{5(i-1)+j+1}=a_{i,j}=e_{i,j}-e_{i,j+1}
 \quad(1\leq i\leq4,\ 0\leq j\leq4),
\]
followed by
\[
 g_{21}=e_1-e_2,\quad g_{22}=e_2-e_3,\quad
 g_{23}=e_3-e_4,\quad g_{24}=e_3+e_4
\]
in the original \(D_4\) coordinates.  All unspecified components
are zero.  Adjoin the five actual glue lifts
\begin{align*}
 g_{25}&=(3\lambda,3\lambda,3\lambda,3\lambda;0),&
 g_{26}&=(3\lambda,3\lambda,0,0;v_D),\\
 g_{27}&=(3\lambda,0,3\lambda,0;s_D),&
 g_{28}&=(2\lambda,2\lambda,2\lambda,0;0),\\
 g_{29}&=(2\lambda,4\lambda,0,2\lambda;0).
\end{align*}
Then
\begin{equation}
 N=\sum_{i=1}^{29}\mathbb Zg_i,\qquad
 H_i=(\rho,g_i)=
 \begin{cases}
 1,&1\leq i\leq24,\\
 30,18,18,15,20,&i=25,26,27,28,29\text{ respectively}.
 \end{cases}
 \label{nca:eq:N-generators}
\end{equation}
Indeed the first 24 vectors span \(R\), the last five generate
the original glue, and subtraction of the corresponding glue lift
leaves every element of \(N\) in \(R\).  The height calculation
uses \((\rho_A,\lambda)=5/2\),
\((\rho_D,v_D)=(\rho_D,s_D)=3\), and the simple-root heights one.

Define the following 30, fully specified neighbour generators:
\begin{equation}
 h_1=6\alpha,\qquad h_i=g_i-H_i\alpha\quad(2\leq i\leq29),
 \qquad h_{30}=v=\rho/6-\alpha.
 \label{nca:eq:Lambda-generators}
\end{equation}
Then \(K=\sum_{i=1}^{29}\mathbb Zh_i\) and
\(\Lambda=\sum_{i=1}^{30}\mathbb Zh_i\).
For if \(X=\sum_i n_ig_i\in K\), then
\[
 X=\sum_{i=2}^{29}n_i(g_i-H_i\alpha)
                         +\left(\sum_{i=1}^{29}n_iH_i\right)\alpha,
\]
and the last coefficient is divisible by six. Conversely every
displayed generator before \(h_{30}\) has height divisible by six.
Adjoining \(h_{30}\) is exactly the original neighbour definition.

\subsection{A complete finite calculation of the cubic value modules}

For a lattice \(L\) and a homogeneous real cubic \(F\), define
its additive value module by
\[
 \mathcal V_F(L)=\sum_{X\in L}\mathbb ZF(X)\subset\mathbb R.
\]
This definition retains every value and takes its generated additive
group. It does not claim that every element of the group is one
individual value of \(F\).

\begin{theorem}[Finite cubic-value generators, including negative integers]
\label{nca:thm:finite-values}
If \(L=\sum_{i=1}^m\mathbb Zb_i\), allowing dependent generators,
then \(\mathcal V_F(L)\) is generated by the following finite list:
\begin{align}
 A_i={}&F(b_i),&&1\leq i\leq m,\label{nca:eq:single}\\
 D^+_{ij}={}&F(b_i+b_j)-F(b_i)-F(b_j),&&i<j,\label{nca:eq:Dplus}\\
 D^-_{ij}={}&F(b_i-b_j)-F(b_i)+F(b_j),&&i<j,\label{nca:eq:Dminus}\\
 D_{ijk}={}&F(b_i+b_j+b_k)-F(b_i+b_j)-F(b_i+b_k)\notag\\
          &{}-F(b_j+b_k)+F(b_i)+F(b_j)+F(b_k),&&i<j<k.
 \label{nca:eq:Dthree}
\end{align}
Consequently \eqref{nca:eq:N-generators} gives exactly 4495
specified generators of \(\mathcal V_F(N)\), and
\eqref{nca:eq:Lambda-generators} gives exactly 4960 specified
generators of \(\mathcal V_F(\Lambda)\), for every fixed
orthonormal octonion frame and its actual \(\phi_{ijk}\).
\end{theorem}

\begin{proof}
Let \(B_F\) be the symmetric trilinear polarization, with
\(B_F(X,X,X)=F(X)\); explicitly its value is
\(1/6\) times the seven-term expression in
\eqref{nca:eq:Dthree}, with its three inputs replacing the three
generators. For every integer tuple \(n\), expansion gives
\[
 F\left(\sum_i n_ib_i\right)
 =\sum_i n_i^3A_i
 +\sum_{i<j}\bigl(n_i^2n_j C_{ij}+n_in_j^2 E_{ij}\bigr)
 +\sum_{i<j<k}n_in_jn_kD_{ijk},
\]
where \(C_{ij}=3B_F(b_i,b_i,b_j)\) and
\(E_{ij}=3B_F(b_i,b_j,b_j)\).  The two pair generators are
\(D^+_{ij}=C_{ij}+E_{ij}\) and
\(D^-_{ij}=-C_{ij}+E_{ij}\). Therefore their contribution is
\[
 \frac{n_in_j(n_i+n_j)}2D^+_{ij}
       +\frac{n_in_j(n_j-n_i)}2D^-_{ij}.
\]
Both coefficients are integers, for all positive, zero, and negative
integers \(n_i,n_j\): if either integer is even the claim is immediate,
and if both are odd their sum and difference are even. This proves
that the finite list generates all values. Conversely each entry
of the list is an integer linear combination of actual values on
\(L\). The proof does not use linear independence of the generators.
Finally \(29+2\binom{29}{2}+\binom{29}{3}=4495\) and
\(30+2\binom{30}{2}+\binom{30}{3}=4960\).
\end{proof}

For numerical arithmetic now specify the additional octonion frame
choice completely.  Use \(\mathbb O=\mathbb H\oplus\mathbb H\ell\)
in ordered coordinates \((a,b)\), with
\begin{equation}
 \begin{gathered}
 (a,b)(c,d)=(ac-\overline d\,b,\ da+b\overline c),\\
 u_0=(1,0),\quad u_1=(i,0),\quad u_2=(j,0),\quad u_3=(k,0),\\
 u_4=(0,1),\quad u_5=(0,i),\quad u_6=(0,j),\quad u_7=(0,k).
 \end{gathered}
 \label{nca:eq:Cayley-frame}
\end{equation}
Here \(ij=k\), \(jk=i\), \(ki=j\), and the reversed products
are negative.  The positive oriented three-form triples are exactly
\begin{equation}
 (123),\ (145),\ (176),\ (246),\ (257),\ (347),\ (365).
 \label{nca:eq:seven-triples}
\end{equation}
Every \(\phi_{ijk}\) is obtained from these by alternation, or is
zero.  This is a specified frame for the retained, previously
frame-dependent isometry, without a change to any lattice coordinates
or to the Albert matrix convention.

\begin{theorem}[Exact value modules and their quotient of order twelve]
\label{nca:thm:value-modules}
For the frame \eqref{nca:eq:Cayley-frame}, the exact additive modules are
\begin{align}
 \mathcal V_F(N)
    &=\frac1{27}\mathbb Z\ \oplus\
       \frac{\sqrt2}{54}\mathbb Z\ \oplus\
       \frac{\sqrt3}{18}\mathbb Z,\label{nca:eq:N-values}\\
 \mathcal V_F(\Lambda)
    &=\frac1{27}\mathbb Z\ \oplus\
       \frac{\sqrt2}{216}\mathbb Z\ \oplus\
       \frac{\sqrt3}{54}\mathbb Z.\label{nca:eq:Lambda-values}
\end{align}
In particular the inclusion of these two additive subgroups of
\(\mathbb R\) has the exact quotient
\begin{equation}
 \mathcal V_F(\Lambda)/\mathcal V_F(N)
                  \cong\mathbb Z/4\oplus\mathbb Z/3
                  \cong\mathbb Z/12.
 \label{nca:eq:value-quotient}
\end{equation}
This is an inclusion of cubic value modules; it does not assert the
false lattice inclusion \(N\subset\Lambda\).
\end{theorem}

\begin{proof}
The polynomial coefficients in
\eqref{nca:eq:P0}--\eqref{nca:eq:P3} are rational on both lattices
for this exact multiplication table.  The numbers
\(1,\sqrt2,\sqrt3\) are linearly independent over \(\mathbb Q\).
For example, isolating one radical in a rational relation and
squaring either gives a rational multiple of another irrational
radical, or a rational square equal to two, three, or \(3/2\),
which is impossible by parity of prime exponents. Thus a value has
a unique rational coefficient triple \((P_0,P_2,P_3)\).

Apply the complete finite list in
Theorem~\ref{nca:thm:finite-values} to the displayed original
generators. Exact substitution gives the following least common
multiples of coefficient denominators. The denominators in each
entry are ordered as \((P_0,P_2,P_3)\), and a zero coefficient has
denominator one:
\begin{equation}
\begin{array}{c|c|c|c|c}
 &A_i&D^+_{ij}&D^-_{ij}&D_{ijk}\\\hline
 N&(27,27,1)&(9,18,18)&(9,18,18)&(9,18,18)\\
 \Lambda&(27,216,27)&(27,108,54)&(27,108,54)&(9,36,54)
\end{array}.
\label{nca:eq:denominator-certificate}
\end{equation}
This table is a finite rational calculation, with no bound on
lattice points or heuristic sampling.  The complete arithmetic
certificate is the accompanying \texttt{verify\_cubic.py}: its
\texttt{residual\_data} and \texttt{cubic} functions are exactly
\eqref{nca:eq:original-s}--\eqref{nca:eq:P3}; it forms every
\(A_i,D^+_{ij},D^-_{ij},D_{ijk}\) for the 29 or 30 explicitly
defined vectors, using rational arithmetic.  The \texttt{cross}
function executes the three cyclic positive and three negative
entries for each of the seven triples
\eqref{nca:eq:seven-triples}.  Hence the full finite calculation
is specified by these equations, the finite index bounds in
Theorem~\ref{nca:thm:finite-values}, and integer arithmetic.
Its 4495 and 4960 entries verify the upper inclusions in
\eqref{nca:eq:N-values}--\eqref{nca:eq:Lambda-values}.

Here are short certificates for the reverse inclusions. Coordinates
in the next two tables are relative to the respective three
proposed generators on the right sides of
\eqref{nca:eq:N-values} and \eqref{nca:eq:Lambda-values}.
In the first table all pair symbols use the \(g_i\):
\begin{equation}
\begin{array}{c|r r r}
 F(g_1)&-4&4&0\\
 F(g_2)&-5&-4&0\\
 F(g_7)&0&18&0\\
 F(g_{28})&-80&52&0\\
 D^+_{21,28}&0&-36&16\\
 D^+_{2,27}&-27&27&3\\
 D^+_{2,25}&-81&45&0
\end{array}.
\label{nca:eq:N-lower-certificate}
\end{equation}
The determinants of the three columns obtained by transposing
rows \((1,4,5)\), \((2,3,6)\), and \((2,7,6)\) are,
respectively, \(1792,-270,-1647\). Their greatest common divisor
is one. In the second table all pair symbols use the \(h_i\):
\begin{equation}
\begin{array}{c|r r r}
 F(h_2)&14&-272&0\\
 F(h_4)&5&-128&0\\
 F(h_{28})&13780&-74672&0\\
 F(h_{30})&14&-63&-2\\
 D^+_{2,27}&7425&-60624&-9
\end{array}.
\label{nca:eq:Lambda-lower-certificate}
\end{equation}
Rows \((1,3,4)\), \((1,4,5)\), and \((2,4,5)\) give
determinants \(-5405504,2315394,1281267\), whose greatest
common divisor is also one.

For clarity, the elementary integer argument using these
determinants does not require a normal-form theorem.  Let
\(W\subset\mathbb Z^3\) be the group generated by one displayed
table.  For any three of its columns forming a matrix \(C\),
the adjugate identity \(C\operatorname{adj}(C)=\det(C)I\)
implies \(\det(C)\mathbb Z^3\subset W\).  Bezout's identity
for the three displayed determinants, whose gcd is one, therefore
gives \(\mathbb Z^3\subset W\). The opposite inclusion is
immediate from the integer entries. Every table entry is an
integer combination of actual values by its definition, proving
the two reverse inclusions.  Finally the coordinates of the
inclusion of the first module in the second are
\((a,b,c)\mapsto(a,4b,3c)\). Its cokernel is the displayed
product of cyclic groups, with order twelve. This also proves
the asserted exact quotient map.
\end{proof}

\subsection{Frame dependence and exact symmetry transport}

The finite-list theorem is valid for the actual coefficients of
every orthonormal frame. The numerical modules
\eqref{nca:eq:N-values}--\eqref{nca:eq:Lambda-values} refer to
the frame \eqref{nca:eq:Cayley-frame}.  This dependence has an
exact comparison map. If \(u_a'=Tu_a\) for a real isometry
\(T\) fixing \(1\), then
\[
 \mathcal L_{u'}=(T,T,T)\mathcal L_u.
\]
The symmetric part \(S\) of the trilinear tensor is unchanged,
because \(T\) preserves the scalar coordinate and inner product.
Writing \(\mathbf m=\operatorname{Im}M\), its exact change is
\begin{equation}
 F_{u'}(X)-F_u(X)
  =-3\bigl(\phi(T\mathbf a,T\mathbf b,T\mathbf m)
                  -\phi(\mathbf a,\mathbf b,\mathbf m)\bigr).
 \label{nca:eq:frame-transport}
\end{equation}
This follows from the complete expansion in the polynomial proof.
If \(T\) is an octonion automorphism it preserves \(\phi\), so
the cubic and its value modules are unchanged.

The dependence on a general orthonormal frame is substantive. Take
the actual lattice vector with the only nonzero fibre coordinates
\[
 x_{1,1}=1,\ x_{1,5}=-1,\quad
 x_{2,2}=1,\ x_{2,4}=-1,\qquad d=(1,1,0,0).
\]
It lies in \(R\subset N\), has \(A=u_1\), \(B=u_2\), and
\(M=(u_4+u_5)/\sqrt3\). In the Cayley frame its value is zero.
Replace only
\[
 u'_3=(2u_3-\sqrt5u_4)/3,\qquad
 u'_4=(\sqrt5u_3+2u_4)/3.
\]
This is an exact orthonormal frame change; the same original
lattice vector now has value \(-\sqrt{15}/3\), because
\(\phi(u_1,u_2,u'_4)=\sqrt5/3\) and
\(\phi(u_1,u_2,u_5)=0\). This value is outside
\(\mathbb Q+\sqrt2\mathbb Q+\sqrt3\mathbb Q\). Indeed the
independent sign changes of \(\sqrt2,\sqrt3\) are automorphisms
of \(\mathbb Q(\sqrt2,\sqrt3)\). If an element \(z\) in that
field had \(z^2=15\), each sign change would send it to
\(z\) or \(-z\). In the basis
\(1,\sqrt2,\sqrt3,\sqrt6\) this makes \(z\) a rational
multiple of a single basis vector. Its square would force one
of \(15,15/2,5,5/2\) to be a rational square, impossible by
parity of its prime exponents.
Thus the numerical specialization is not silently extended to
an unspecified orthonormal frame.

Return to the explicit Cayley frame for the following numerical
witnesses. Let \(\sigma=\rho(289)\) act simultaneously by
\(e_{i,j}\mapsto e_{i,j+1\bmod6}\), and trivially on \(D_4\).
The exact isometry gives
\(\mathcal L\sigma^2=c\mathcal L\) for
\(c(q_1,q_2,q_3)=(q_3,q_1,q_2)\). Cyclic real associativity
therefore proves \(F(\sigma^{2k}X)=F(X)\) for every real \(X\).
These maps preserve \(N\) by their original discriminant action.
The stabilizer of this cubic within the six-element residue group
is exactly its three-element subgroup: direct substitution at
\(\alpha=g_1\) gives
\begin{equation}
 F(\alpha)=\frac{-4+2\sqrt2}{27},\qquad
 F(\sigma\alpha)=\frac{-5-2\sqrt2}{27},
 \label{nca:eq:residue-nonpreserver}
\end{equation}
which differ. All odd powers have the same effect on the cubic
as \(\sigma\), by its proved invariance under even powers.

The original reflection exchanging \(e_{1,0},e_{1,1}\) sends
\(\alpha\) to \(-\alpha\), hence changes its nonzero cubic
value to its negative. The original full glue-preserving lifts
\[
 \begin{aligned}
 \widehat R(X_1,X_2,X_3,X_4;d)
       &=(X_2,X_1,X_3,-X_4;R_Dd),\\
 \widehat H(X_1,X_2,X_3,X_4;d)
       &=(X_1,X_3,X_2,-X_4;H_Dd).
 \end{aligned}
\]
also have explicit witnesses:
\[
 F(\widehat R\alpha)=0,\qquad
 F(g_7)=\frac{\sqrt2}{3},\qquad F(\widehat Hg_7)=0.
\]
Their original matrices remain
\[
 R_D=\operatorname{diag}(-1,1,1,1),\qquad
 H_D=\frac12\begin{pmatrix}1&1&1&1\\1&1&-1&-1\\
                           1&-1&1&-1\\1&-1&-1&1\end{pmatrix}.
\]
These maps preserve the original lattice, as can be checked on its
displayed generators.  Both \(R_D\) and \(H_D\) preserve \(D_4\):
the former changes one sign, while the latter has integral coordinates
on every even-sum integer vector and has output sum twice its first
input coordinate. They are orthogonal involutions. Their marked
discriminant actions fix \(v_D\) and exchange \(s_D,c_D\), and
exchange \(v_D,s_D\) while fixing \(c_D\), respectively.
On the binary glue generators, \(\widehat R\) fixes the first
two and sends the third to their second-plus-third sum;
\(\widehat H\) fixes the first and exchanges the second and third.
On the ternary code generators \(z_1=(1,1,1,0)\),
\(z_2=(1,2,0,1)\), the two signed component permutations act as
\((z_1,z_2)\mapsto(z_1,2z_2)\) and
\((z_1,z_2)\mapsto(z_1,2z_1+2z_2)\), respectively.
Thus both preserve the full original glue and its inverse image.
Every within-fibre coordinate permutation also preserves that inverse
image, since its action on \(\lambda+A_5\) is trivial:
\(Q\lambda-\lambda=Qe_0-e_0\in A_5\).
These are concrete failures to preserve this cubic, without
discarding their actual lattice actions or glue.

Negation preserves each lattice \(N\) and \(\Lambda\), and its
exact polynomial law is \(F(-X)=-F(X)\). In the neighbour the
nonzero witness is its displayed norm-four generator:
\begin{equation}
 F(\rho)=15\sqrt2+16\sqrt3,\qquad
 F(v)=\frac{14}{27}-\frac{7\sqrt2}{24}-\frac{\sqrt3}{27}.
 \label{nca:eq:special-values}
\end{equation}
Both equalities follow by substituting the original coordinates
of \(\rho\) and \(v\) into the proved polynomial.

Every original lattice isometry \(T\in O(N)\), whether or not
it preserves the cubic, has the exact replacement transport
\begin{equation}
 \Gamma_T=\mathcal LT\mathcal L^{-1},\qquad
 F^T(Y)=F(T^{-1}Y),\qquad F^T(TX)=F(X),
 \label{nca:eq:all-transport}
\end{equation}
and
\[
 T\Lambda_\rho=\Lambda_{T\rho},\qquad
 \mathcal V_{F^T}(\Lambda_{T\rho})=\mathcal V_F(\Lambda_\rho),
 \qquad \mathcal V_{F^T}(N)=\mathcal V_F(N).
\]
For the neighbour transport, write its exact elements as
\(x+k\rho/6\), where \(x\in N\) and
\((\rho,x)+k\equiv0\pmod6\); application of \(T\) gives
the identical congruence with \((T\rho,Tx)\). The inverse is
\(T^{-1}\). The value-module equalities then follow by the
bijections and the displayed identity \(F^T(TX)=F(X)\).
In particular for \(T=\sigma^{2k}\), one has \(F^T=F\), so
the same retained cubic has exactly the same value module on
every neighbour \(\Lambda_{\sigma^{2k}\rho}\) in this cyclic
orbit. This statement names its target neighbour and does not
assume the chosen \(\rho\) is fixed.
